\documentclass{article}

\usepackage[margin=1in]{geometry}
\usepackage{graphicx, amsmath, amsthm, amssymb, enumitem, lipsum, hyperref}

\setlist[enumerate]{leftmargin=1.25cm}

\newcommand{\N}{\mathbb{N}}
\newcommand{\Z}{\mathbb{Z}}
\newcommand{\Q}{\mathbb{Q}}
\newcommand{\R}{\mathbb{R}}
\newcommand{\C}{\mathbb{C}}

\newcommand{\rotvert}{\rotatebox[origin=c]{90}{$\vert$}}

\newenvironment{note}[2][Note]{\begin{trivlist}
\item[\hskip \labelsep {\bfseries #1}\hskip \labelsep {\bfseries #2.}]}{\end{trivlist}}

\newenvironment{fact}[2][Fact]{\begin{trivlist}
\item[\hskip \labelsep {\bfseries #1}\hskip \labelsep {\bfseries #2.}]}{\end{trivlist}}

\newenvironment{definition}[2][Definition]{\begin{trivlist}
\item[\hskip \labelsep {\bfseries #1}\hskip \labelsep {\bfseries #2.}]}{\end{trivlist}}

\newenvironment{proposition}[2][Proposition]{\begin{trivlist}
\item[\hskip \labelsep {\bfseries #1}\hskip \labelsep {\bfseries #2.}]}{\end{trivlist}}

\newenvironment{principle}[2][Principle]{\begin{trivlist}
\item[\hskip \labelsep {\bfseries #1}\hskip \labelsep {\bfseries #2.}]}{\end{trivlist}}

\newenvironment{lemma}[2][Lemma]{\begin{trivlist}
\item[\hskip \labelsep {\bfseries #1}\hskip \labelsep {\bfseries #2.}]}{\end{trivlist}}

\newenvironment{theorem}[2][Theorem]{\begin{trivlist}
\item[\hskip \labelsep {\bfseries #1}\hskip \labelsep {\bfseries #2.}]}{\end{trivlist}}

\newenvironment{corollary}[2][Corollary]{\begin{trivlist}
\item[\hskip \labelsep {\bfseries #1}\hskip \labelsep {\bfseries #2.}]}{\end{trivlist}}

\newenvironment{envsection}[1]{\begin{trivlist}
\item[\hskip \labelsep {\bfseries #1}]}{\end{trivlist}}


\begin{document}
\setlength{\abovedisplayskip}{4pt}
\setlength{\belowdisplayskip}{4pt}


\begin{center}
    \Large{\textbf{Homework 5}}\\
    \normalsize March 12 - 19
\end{center}
\vspace{10pt}

\subsection*{Reading}

\begin{envsection}{1. Speech and Language Processing:}
    I added this book to the GitHub. Start reading Chapter 3: $N$-gram Language Models. This section will put MLE in the context of natural language processing, so it should feel pretty natural after our linear regression discussion.
\end{envsection}

\subsection*{Problems}

\begin{envsection}{Problem 1.}
    Use induction to prove the \textit{chain rule of probability}: For $n$ random variables $X_1,\dots,X_n$,
    \[
    P(X_1,\dots,X_n) = \prod_{k=1}^{n} P(X_k \mid X_{1:k-1}).
    \]
    (Here, $X_{1:k-1}$ refers to the joint distribution $X_1,\dots,X_{k-1}$).
\end{envsection}


\end{document}